\documentclass[11pt]{article}

\usepackage[margin=1in]{geometry}
\usepackage{hyperref}
\usepackage{amsmath}
\usepackage{enumitem}
\usepackage{listings}
\usepackage{xcolor}

\hypersetup{
    colorlinks=true,
    linkcolor=blue,
    urlcolor=blue,
    citecolor=blue
}

\title{CS 300 -- Week 5 Resources\\\large Binary Search Trees, Traversals, and C++ File I/O}
\author{Jose de Lima}
\date{\today}

\begin{document}
\maketitle

\section*{Overview}
Week 5 brings together three threads. First, the \textbf{binary search
tree (BST)} -- a hierarchical data structure that keeps keys in sorted
order and supports $O(\log n)$ operations on average. Second, the
standard \textbf{tree traversals} (in-order, pre-order, post-order) that
we use to visit every node in a well-defined sequence. Third, practical
\textbf{C++} foundations and \textbf{file I/O}, since the CS 300 projects
read bid data from CSV files.

\section{Binary Search Trees}
\textbf{Resource:} \href{https://www.youtube.com/watch?v=kq0LPPNwvzQ}{CS 260 Lesson 6 -- Binary Search Tree (YouTube)}

\subsection*{Definition}
A \textbf{binary search tree} is a binary tree in which every node
satisfies the \textbf{BST property}:
\begin{itemize}[leftmargin=*]
    \item Every key in the node's \emph{left} subtree is \emph{less than}
          the node's key.
    \item Every key in the node's \emph{right} subtree is \emph{greater
          than} the node's key.
    \item Both subtrees are themselves BSTs.
\end{itemize}
That one invariant is what lets us throw away half of the remaining
search space at every step, giving logarithmic lookups when the tree
stays balanced.

\subsection*{Node Layout}
\begin{lstlisting}[basicstyle=\ttfamily\small,frame=single]
struct Node {
    Key   key;
    Value value;
    Node* left;
    Node* right;
};
\end{lstlisting}

\subsection*{Core Operations}

\paragraph{Search.}
Compare the target key to the current node's key. Go left if smaller,
right if larger, stop if equal. Return ``not found'' if you fall off a
leaf.

\paragraph{Insert.}
Search for the key; when you reach a \texttt{nullptr} child, attach a
new node there. Duplicates are usually either rejected or treated as
an update.

\paragraph{Remove.}
Three cases once the node is found:
\begin{enumerate}[leftmargin=*]
    \item \textbf{Leaf} -- free it and set the parent's pointer to
          \texttt{nullptr}.
    \item \textbf{One child} -- splice the child up into the node's
          slot.
    \item \textbf{Two children} -- replace the node's key/value with
          its \textbf{in-order successor} (smallest key in the right
          subtree), then delete that successor from the right subtree.
\end{enumerate}

\subsection*{Pseudocode (insert \& search)}
\begin{lstlisting}[basicstyle=\ttfamily\small,frame=single]
function search(node, key):
    if node == nullptr: return NOT_FOUND
    if key == node.key: return node.value
    if key <  node.key: return search(node.left,  key)
    else:               return search(node.right, key)

function insert(node, key, value):
    if node == nullptr: return new Node(key, value)
    if key <  node.key: node.left  = insert(node.left,  key, value)
    elif key > node.key: node.right = insert(node.right, key, value)
    else: node.value = value          # upsert
    return node
\end{lstlisting}

\subsection*{Complexity}
\begin{itemize}[leftmargin=*]
    \item \textbf{Balanced tree:} insert, search, remove are
          $O(\log n)$.
    \item \textbf{Worst case (degenerate / sorted input):} the tree
          turns into a linked list, giving $O(n)$. Self-balancing
          variants (AVL, red-black) avoid this.
    \item \textbf{Space:} $O(n)$ for the nodes, plus $O(h)$ recursion
          stack for traversals, where $h$ is the height.
\end{itemize}

\section{BST Traversals}

A \textbf{traversal} is a rule for visiting every node in the tree
exactly once. For BSTs the three depth-first traversals differ only in
\emph{when} the current node is visited relative to its children.

\subsection*{In-Order (Left, Node, Right)}
\begin{lstlisting}[basicstyle=\ttfamily\small,frame=single]
function inorder(node):
    if node == nullptr: return
    inorder(node.left)
    visit(node)
    inorder(node.right)
\end{lstlisting}
In a BST, in-order traversal produces keys in \textbf{sorted ascending
order}. This is the traversal the Project Two bid-list printout uses.

\subsection*{Pre-Order (Node, Left, Right)}
\begin{lstlisting}[basicstyle=\ttfamily\small,frame=single]
function preorder(node):
    if node == nullptr: return
    visit(node)
    preorder(node.left)
    preorder(node.right)
\end{lstlisting}
Visits the root before its subtrees. Useful for \textbf{copying} or
\textbf{serializing} a tree, because the root is emitted first and can
be reconstructed top-down.

\subsection*{Post-Order (Left, Right, Node)}
\begin{lstlisting}[basicstyle=\ttfamily\small,frame=single]
function postorder(node):
    if node == nullptr: return
    postorder(node.left)
    postorder(node.right)
    visit(node)
\end{lstlisting}
Visits children before the root. Useful for \textbf{deleting} or
\textbf{freeing} a tree -- you must not free a parent before its
children, or you lose the pointers you need.

\subsection*{Example}
\begin{center}
\begin{tabular}{c}
\qquad\qquad 8 \\
\qquad\quad / \quad \textbackslash \\
\qquad 3 \qquad 10 \\
\qquad / \textbackslash \qquad \textbackslash \\
\quad 1 \quad 6 \qquad 14
\end{tabular}
\end{center}
\begin{itemize}[leftmargin=*]
    \item In-order:\quad   1, 3, 6, 8, 10, 14
    \item Pre-order:\quad  8, 3, 1, 6, 10, 14
    \item Post-order:\quad 1, 6, 3, 14, 10, 8
\end{itemize}

\subsection*{Complexity}
All three traversals are $O(n)$ time (every node is visited once) and
$O(h)$ additional space for the recursion stack.

\subsection*{Breadth-First (level-order) -- Honorable Mention}
Not covered in the PDF, but worth knowing: a queue-driven traversal
that visits every node on depth 0, then depth 1, and so on -- useful
for shortest-path-by-edges problems and for printing the tree
``layer by layer.''

\section{C++ Foundations}
\textbf{Resource:} \href{https://www.w3schools.com/cpp/default.asp}{W3Schools C++ Tutorial}

The W3Schools tutorial is a fast, example-driven refresher on the parts
of C++ the CS 300 projects lean on. The sections most worth reviewing
for this course are:
\begin{itemize}[leftmargin=*]
    \item \textbf{Syntax, variables, and data types} -- including
          \texttt{int}, \texttt{double}, \texttt{bool},
          \texttt{std::string}.
    \item \textbf{Control flow} -- \texttt{if}/\texttt{else},
          \texttt{switch}, \texttt{while}, \texttt{for},
          range-based \texttt{for}.
    \item \textbf{Functions and references} -- pass-by-value vs.\
          pass-by-reference, default arguments, overloading.
    \item \textbf{Classes and objects} -- constructors, destructors,
          \texttt{this}, access modifiers. Every project structure
          (\texttt{Bid}, \texttt{Node}, \texttt{HashTable},
          \texttt{BinarySearchTree}) is a class.
    \item \textbf{Pointers, references, and dynamic memory} --
          \texttt{new} / \texttt{delete}, \texttt{nullptr},
          pointer arithmetic. This is where linked lists, hash
          table chains, and BST nodes live.
    \item \textbf{The STL containers} -- \texttt{std::vector},
          \texttt{std::string}. They back the course's starter code.
\end{itemize}

\subsection*{Things to Lock In for CS 300}
\begin{itemize}[leftmargin=*]
    \item The difference between an \textbf{object}, a
          \textbf{pointer to an object}, and a \textbf{reference} to
          one -- and when to use each.
    \item How to write a correct \textbf{destructor} that frees every
          node in a linked structure (use post-order recursion).
    \item How to pass a \texttt{std::vector} without copying it --
          use \texttt{const std::vector<T>\&}.
\end{itemize}

\section{C++ File I/O}
\textbf{Resource:} \href{https://www.w3schools.com/cpp/cpp_files.asp}{W3Schools -- C++ Files}

The CS 300 projects read a CSV file of bids into memory, so file I/O
is not optional. C++ exposes three stream classes in \texttt{<fstream>}:
\begin{itemize}[leftmargin=*]
    \item \texttt{std::ofstream} -- write-only output file stream.
    \item \texttt{std::ifstream} -- read-only input file stream.
    \item \texttt{std::fstream}  -- read/write file stream.
\end{itemize}

\subsection*{Reading a File Line-by-Line}
\begin{lstlisting}[basicstyle=\ttfamily\small,frame=single]
#include <fstream>
#include <string>
#include <iostream>

std::ifstream in("bids.csv");
if (!in.is_open()) {
    std::cerr << "Could not open bids.csv\n";
    return 1;
}

std::string line;
while (std::getline(in, line)) {
    // parse `line` into a Bid and insert into your data structure
}
in.close();
\end{lstlisting}

\subsection*{Key Points}
\begin{itemize}[leftmargin=*]
    \item Always check \texttt{is\_open()} (or the stream's boolean
          conversion) before reading -- a missing file does not throw.
    \item \texttt{std::getline} reads until the next newline; pair it
          with \texttt{std::stringstream} to split a CSV row on
          commas.
    \item RAII: when the \texttt{ifstream} goes out of scope the file
          is closed automatically. Calling \texttt{close()} is
          optional but makes intent explicit.
    \item Use \texttt{std::ofstream} in \texttt{std::ios::app} mode to
          \emph{append} instead of overwrite.
\end{itemize}

\subsection*{Parsing a CSV Row}
\begin{lstlisting}[basicstyle=\ttfamily\small,frame=single]
#include <sstream>

std::stringstream ss(line);
std::string field;
std::vector<std::string> cells;
while (std::getline(ss, field, ',')) {
    cells.push_back(field);
}
// cells[0], cells[1], ... are the columns
\end{lstlisting}

\section*{Key Takeaways}
\begin{itemize}[leftmargin=*]
    \item A BST's power comes from one invariant -- left $<$ node $<$
          right -- which enables $O(\log n)$ operations when the tree
          stays balanced.
    \item In-order, pre-order, and post-order differ only in when the
          node is visited: use in-order to \emph{list} keys in order,
          pre-order to \emph{copy/serialize}, post-order to
          \emph{free/destroy}.
    \item Brush up on C++ pointers, classes, and destructors before
          implementing the project -- most bugs in this course come
          from mishandled memory, not the algorithms themselves.
    \item \texttt{std::ifstream} + \texttt{std::getline} +
          \texttt{std::stringstream} is the standard recipe for
          loading a CSV in C++.
\end{itemize}

\section*{Links Summary}
\begin{itemize}[leftmargin=*]
    \item Binary Search Tree video: \url{https://youtu.be/kq0LPPNwvzQ}
    \item W3Schools -- C++ Tutorial: \url{https://www.w3schools.com/cpp/default.asp}
    \item W3Schools -- C++ Files: \url{https://www.w3schools.com/cpp/cpp_files.asp}
\end{itemize}

\end{document}
